% Nouveau EGA V (cote 25): archive cote 25, site page 80 (PDF page 81).
% typescript. Draft compared: none.
% AI-transcribed by Claude Opus 5.5 (claude-opus-5-5, adaptive thinking, effort high), Message Batches API, batch ega25b (msgbatch_01A7n5YuXubUPNK1zsepi8RG), 2026-11-13; checked against both supplied scan parts; see CHECK_DECISIONS.tsv.
\scanpage{25/80}
\mspage{\hand{75}}

\placedmargin{|}m, en notant qu'à changement de corps de base près, $\underline{X}^{(m)}_{\hand{\unsure{\eta|\xi}}}$ s'obtient en
prenant une L' de codimension m-1 passant par x, une H de \unsure{codimension|codimemsion}
1 passant par x, L' et H étant génériques indépendants pour ces proprié
tés (i.e. en termes orthodoxes, on se place au point générique du 
schéma des couples (L',H)  ) et en prenant la section linéaire
de X par $L'\hand{\cap} H$ : on peut commencer par prendre a section par H, qui
est lisse par ce qui précède, puis par L', qui est lisse par l'hypothèse
de récurrence. Ce genre de raisonnement, déjà utilisé pour généraliser
2.6. p.ex. aux sections linéaires de codimension m queconque, mérite
d'être explicité une bonne fois en général, pour pouvoir y référer sans
entrer  chaque fois dans les détails un peu lourds d'une présentati-%
on en forme.

Reste à prouver le corollaire annoncé de 8.8. dans le cas m=1 .
Lorsque  codim T $\hand{\geqslant}$ 2, comme d'autre part l'hyperplan Q de $\overset{\hand{\vee}}{P}$ des
$H_{\hand{\xi}}$ tels que $x\hand{\in} H_{\hand{\xi}}$ est de codimension 1, son point générique $\eta$ ne peut
être élement de T, et on gagne. Dans le cas codim T = 1, comme T est
irréductible, on ne peut avoir $\hand{\eta}\hand{\in} T$ que si Q=T ,
i.e. en termes géométriques  (supposant k alg clos, ce qui est loisible) pour tout $z\hand{\in} X(k)$ l'espace tangent
de X en z (ou plutôt son image par $f'_{z}$ ) passe par x. Montrons que
ceci ne peut se présenter lorsque $\underline{X}^{\mathrm{sing}}\hand{\longrightarrow}T$ est gén étale, i.e.
lorsqu'on est sous les conditions de 8.8., sauf si f(X) = \hand{\{}x\hand{\}} donc X
de dimension 0. En effet, 8.7. c) (qui exprime essentiellement la
symétrie  dans la relation entre X et sa \hand{``}duale\hand{''} T\hand{)} implique alors que
pour presque tout point $z\hand{\in} X(k)$, f(z) est orthogonal à l'espace tangent
à T en un certain point, donc (puisque T=Q) orthogonal à Q, d'où f(z)=x, d'où
f(X) =\hand{\{}x\hand{\}} . Cela achève la démonstration de nos remords, donc de 13.2.

\placedmargin{80}
